\section{Introduction}


The principle of Occam's razor---that simpler explanations are preferable---is a foundational concept in machine learning. Algorithmic information theory provides an elegant formalization through the \emph{Minimum Description Length}~\citep[MDL;][]{rissanen1978modeling} principle. The MDL principle frames learning as data compression: the best model for a dataset is the one that minimizes the combined length of the model's description plus the description of the data encoded with that model. Intuitively, any regularity in the data that is useful for prediction is also useful for compression, and vice versa.

The empirical success of large neural networks might seem to contradict the MDL principle if we consider a naive, uniform encoding of their weights. Such networks often generalize surprisingly well, even when highly over-parameterized and trained with simple maximum likelihood objectives~\citep{zhang2017understanding}. However, various methods have been proposed for training compressible neural networks based on, e.g., quantization~\citep{han2016deep}, subspace training~\citep{li2018measuring}, low-rank approximation, variational inference~\citep{louizos2017bayesian,hinton1993keeping}, or some combination~\citep{lotfi2022pac,lotfi2024non,demoss2025complexity}. These methods have highlighted that neural networks can often be highly compressed, i.e. their true complexity can be much smaller than naive parameter counting would suggest~\citep{blier2018description}. While such methods have thus far not led directly to MDL-inspired regularizers that reliably improve generalization, these methods have been successful for establishing tighter compression bounds~\citep{lotfi2022pac,lotfi2024non} and inspiring parameter-efficient fine-tuning methods such as LoRA ~\citep{hu2022lora}.

However, key theoretical questions remain. These methods from prior work can be interpreted as defining different description length measures over a network's weights. Different measures may capture different types of regularities, but fail to capture \emph{all} potential regularities in the network. MDL objectives based on such measures may therefore fail to incentivize capturing all of the regularities in the data, if certain patterns cannot be encoded efficiently in the network's weights. This can lead to sub-optimal compression and---per the MDL principle---sub-optimal generalization. Ideally, we want a description length objective that encourages the model to capture \emph{any} regularity in the data. Assuming perfect optimization, such an objective would lead to optimal compression, regardless of the dataset. To what extent can we implement such an idealized objective?

Algorithmic information theory offers a framework to address this question~\citep{li2008introduction,hutter2005universal,schmidhuber1997discovering}. MDL was inspired by Solomonoff's theory of inductive inference~\citep{solomonoff1964formal}, which proposed to favor hypotheses with low Kolmogorov complexity. Put simply, the Kolmogorov complexity of an object is the length of the shortest program that generates that object (see Section~\ref{sec:background}). Kolmogorov complexity offers optimal compression relative to any other computable description length measure, up to an additive constant. This \emph{universality} is rooted in the Church-Turing thesis that any sufficiently powerful model of computation can simulate any other. Computable resource-bounded approximations maintain this universality in the limit as resource bounds increase. The notion of Kolmogorov complexity can be directly applied to program induction -- priors based on program length have been empirically successful for inducing programs~\cite[e.g.,][]{romera2024mathematical} and grammars~\cite[e.g.,][]{shaw2021compositional} with strong generalization. However, applying this notion to neural networks is less straightforward. While we can easily compute the length of a discrete program, it is less clear how to quantify the complexity of the function a neural network computes in an analogous way. Therefore, a conceptual gap remains between theoretical frameworks based on algorithmic information theory and practical description length objectives for neural networks. We aim to narrow this gap. 

Building on the theory of Kolmogorov complexity and the computational universality of Transformers, we demonstrate the existence of \emph{asymptotically optimal} description length objectives for Transformers, with optimality guarantees analogous to those that hold for Kolmogorov complexity. Our theoretical framework therefore highlights a potential path forward for identifying description length objectives for neural networks that select for models offering greater compression -- and potentially greater generalization. Specifically, we develop the theory of asymptotically optimal description length measures for neural networks through the following contributions:
\begin{itemize}
    \item We define \emph{universal two-part codes} for probabilistic models whose minimum description length for any data sample is provably optimal, up to an additive constant, relative to \emph{any} other two-part code. %
    This framework circumvents the need for arbitrary domain-specific priors by leveraging the universal properties of Kolmogorov complexity. (Section~\ref{sec:two-part-codes})
    \item We prove the existence of \emph{asymptotically optimal families of codes for Transformers}, which are universal in the limit as resource bounds increase. This result is established via a new demonstration that Transformer encoders are computationally universal in their ability to represent any computable, rational-valued conditional probability distribution. (Section~\ref{sec:two-part-codes})
\item We further prove that \emph{tractable and differentiable} objectives for Transformers can be asymptotically optimal, establishing this by constructing and analyzing a variational objective based on an adaptive Gaussian mixture prior. (Section~\ref{sec:variational-codes})
    \item We analyze various codelength objectives \emph{empirically} and analytically. Using a manually constructed solution for an algorithmic task, we show that our variational objective  selects for a highly compressible model with strong generalization. However, we find that standard optimizers fail to discover such low-complexity solutions from a random initialization, highlighting a key challenge for future work. Additionally, we analytically derive asymptotic bounds for various alternative two-part codes. (Section~\ref{sec:experiments})

\end{itemize}

